% Held out of the report (2026-11-26, user): older DRTC results wait for the new robot results (E4).
% Not built; kept unchanged for later.

\section{Setup v2 on the bench}\label{sec:v2bench}

\subsection{\drtc{} against upstream, and the ablations (E1, E2)}

Tables~\ref{tab:v2e1sort} and~\ref{tab:v2e1cube} compare \drtc{} with patient upstream under every fault type: the
stall fraction, how many recoveries succeeded, and how long they took, with the runs that upstream as shipped ended.
Tables~\ref{tab:v2e2sort} and~\ref{tab:v2e2cube} give the ablations: for each part of \drtc{} removed, in how many
policies the stalls rose and by how much. The primary comparisons (\drtc{} against patient upstream, each fault type)
are tested in the plans' reports.

\paragraph{\drtc{} against the patient client (E1).} On sorting (Table~\ref{tab:v2e1sort}), without faults the two
clients stalled about equally: both keep one request in flight, the long-chunk policies rarely run dry, and the
short-chunk ones stall as before. The faults separate them. Under the fault mix \drtc{} cut the stall fraction from
\R{v2/sort_one/e1/mix/act/patient} to \R{v2/sort_one/e1/mix/act/drtc}\,\% for ACT, from \R{v2/sort_one/e1/mix/smolvla/patient} to \R{v2/sort_one/e1/mix/smolvla/drtc}\,\% for
\smolvla{} and from \R{v2/sort_one/e1/mix/pi05/patient} to \R{v2/sort_one/e1/mix/pi05/drtc}\,\% for \pizerofive{}, and it recovered from every
fault for six of the seven policies (GR00T N1.7: \R{v2/sort_one/e1/mix/grootn17/drtcrec}), against \R{v2/sort_one/e1/mix/smolvla/patientrec} for
\smolvla{} and \R{v2/sort_one/e1/mix/groot/patientrec} for \groot{} with the patient client. The difference comes from lost
replies. After an outage the patient client recovered in \R{v2/sort_one/e1/outage/smolvla/patientrec} \smolvla{} runs and
\R{v2/sort_one/e1/outage/groot/patientrec} \groot{} runs, \drtc{} in all of them, \R{v2/sort_one/e1/outage/smolvla/drtcrectime} and
\R{v2/sort_one/e1/outage/groot/drtcrectime}\,s after the outage ended. The patient client sends a fresh request only once, when its
queue runs dry; when that request is lost too, nothing more is sent until LeRobot's 10\,s startup deadline ends the
run. A latency spike alone was absorbed by both clients.

\paragraph{Where \drtc{} was worse: a rising latency.} When every reply became 0.4\,s slower for 5\,s (rise),
\drtc{} stalled more than the patient client for every policy, for example \smolvla{} \R{v2/sort_one/e1/rise/smolvla/patient} against
\R{v2/sort_one/e1/rise/smolvla/drtc}\,\% and ACT \R{v2/sort_one/e1/rise/act/patient} against \R{v2/sort_one/e1/rise/act/drtc}\,\%. The cooldown is the cause. It
is set from the replies seen so far, so a reply that is only slower than the recent ones is given up and a new request
sent; the server, still computing the abandoned one, answers BUSY, and the replacement starts late. Without the
cooldown the stall under rise changed by \R{v2/sort_one/e2/a1_no_cooldown/rise/delta} points (Table~\ref{tab:v2e2sort}). A cooldown
that adapts faster to a slower link is future work (Section~\ref{sec:limits}).

\paragraph{Which part matters (E2).} The cooldown carries almost all of \drtc's benefit. Without it the stall rose in
\R{v2/sort_one/e2/a1_no_cooldown/mix/worse} policies under the fault mix, by a median of \R{v2/sort_one/e2/a1_no_cooldown/mix/delta} points (outage
\R{v2/sort_one/e2/a1_no_cooldown/outage/delta}, drops \R{v2/sort_one/e2/a1_no_cooldown/drops/delta}). The latency estimator made no measurable
difference (median change \R{v2/sort_one/e2/a2_upstream_estimator/mix/delta} points under the mix, \R{v2/sort_one/e2/a2_upstream_estimator/rise/delta}
under rise): on the bench's steady link the maximum of recent turnarounds and the Jacobson--Karels estimate are nearly
the same. The early trigger helped a little (\R{v2/sort_one/e2/a3_upstream_trigger/mix/delta} points under the mix). \drtc's original
cap of the cooldown at half a chunk, which our first fix removed, raised the stall by \R{v2/sort_one/e2/a4_half_chunk_cap/none/delta}
points even without faults and by \R{v2/sort_one/e2/a4_half_chunk_cap/mix/delta} under the mix: the fix holds for all seven policies.

\begin{figure}[htbp]
  \centering
  \genfigure{fig_v2_e1_sort_one}
  \par\medskip
  \genfigure{fig_v2_e1_cubebox_v1}
  \caption{E1: each policy's median stall fraction with patient upstream (orange) and with \drtc{} (green), per
           fault type; sorting (top) and cube in box v1 (bottom). A line going down: \drtc{} stalls less.}
  \label{fig:v2e1}
\end{figure}

\begin{table}[htbp]
  \centering
  \small
  \begin{adjustbox}{max width=\textwidth}
  \gentable{tab_v2_e1_sort_one}
  \end{adjustbox}
  \caption{E1 on sorting: \drtc{} against patient upstream (medians over 8 runs; recovered: runs that moved again
           after each timed fault, summed over the faults of the type).}\label{tab:v2e1sort}
\end{table}

\begin{table}[htbp]
  \centering
  \small
  \begin{adjustbox}{max width=\textwidth}
  \gentable{tab_v2_e1_cubebox_v1}
  \end{adjustbox}
  \caption{E1 on cube in box v1, as Table~\ref{tab:v2e1sort}.}\label{tab:v2e1cube}
\end{table}

\begin{table}[htbp]
  \centering
  \small
  \gentable{tab_v2_e2_sort_one}
  \caption{E2 on sorting: each ablation against \drtc{} (policies whose median stall fraction rose, and the median
           change in percentage points).}\label{tab:v2e2sort}
\end{table}

\begin{table}[htbp]
  \centering
  \small
  \gentable{tab_v2_e2_cubebox_v1}
  \caption{E2 on cube in box v1, as Table~\ref{tab:v2e2sort}.}\label{tab:v2e2cube}
\end{table}

\subsection{The smoothing methods alone (E3a)}

Tables~\ref{tab:v2e3asort} and~\ref{tab:v2e3acube} give each smoothing method alone on \drtc{} under the mix: the jerk
relative to \drtc{} alone, the largest step, the deviation from the policy's plan and the stall fraction. Each method
did one job, and the same one on both tasks.

\paragraph{Two groups of policies.} With \drtc{} alone under the mix, ACT stalled for about 3\% of the run and
\smolvla{} for 12--15\%, while \pizerofive{} (about 38\%), Diffusion Policy (48--53\%), MolmoAct~2 (55--57\%),
\groot{} (64\%) and GR00T~N1.7 (68\%) stalled for a third to two thirds of it: their chunk runs out before the next
one arrives. This split decides which methods help.

\paragraph{The Butterworth filter} lowered the jerk for every policy, to between 0.12 and 0.28 of \drtc{} alone on
sorting, and roughly halved the largest step. Its cost is a deviation of 0.9--1.6\degree{} from the policy's plan; it
did not change the stalls.

\paragraph{Adaptive playback} lowered the stalls of the slow policies, which is what it is for: \groot{} from 64 to
40\%, MolmoAct~2 from 55--57 to 30--35\%, Diffusion Policy from 48--53 to 29--30\% and GR00T~N1.7 from 68 to
53--55\%. It did not help \pizerofive, and ACT and \smolvla{} had little stalling to remove. It also roughly halved
the jerk of the slow policies, because the arm stops and restarts less often, and it never moves the arm off the plan.
Combined with the Butterworth filter (the first run's condition~4) it gave both effects: MolmoAct~2 on cube in box went
from 57\% stalls and a jerk of 232 to 32\% and 15 ($10^3$\degree/s$^3$).

\paragraph{The resume ramp} helps where the arm restarts often: \groot's jerk fell to 0.18 of \drtc{} alone, at a
deviation of 2\degree. \textbf{Zero-lag smoothing} halved the jerk of the fast policies (ACT, \smolvla) but raised it
for \groot{} and MolmoAct~2. \textbf{Temporal ensembling} changed little except for ACT.

\paragraph{RTC} did not help. It more than doubled \smolvla's stalls (from about 14 to 36--38\%): each RTC request
is slower, and the longer turnaround outweighs the smoother joins. For \groot{} the jerk changed by $-19\%$ on sorting
and $+11\%$ on cube in box, within the run-to-run spread.

\paragraph{The stack.} By the rule of Section~\ref{sec:stackrule}, fixed before E3a ran, the stack is \R{v2/stack}.
Both choices were close. The seam choice was decided by one pair (\groot{} on sorting, where RTC lowered the jerk and
ensembling slowed the recovery from the spike); in the other three RTC-capable pairs both candidates were ruled out.
Of the fourteen pairs for the jitter choice, the Butterworth filter won two (ACT on both tasks) and zero-lag smoothing
one; in the other eleven both were ruled out, the filter mostly by the deviation guard (1.0--1.6\degree{} against a
limit of 1\degree) and zero-lag smoothing by slower recovery. Eight runs per cell also leave a wide spread: the same
\drtc{} cell of \groot{} on sorting had a median jerk of 178 in one plan and 219 in the other, so differences below
about 20\% are not reliable.

\begin{table}[htbp]
  \centering
  \small
  \begin{adjustbox}{max width=\textwidth}
  \gentable{tab_v2_e3a_sort_one}
  \end{adjustbox}
  \caption{E3a on sorting: each smoothing method alone on \drtc{} under the fault mix (medians over 8 runs).}\label{tab:v2e3asort}
\end{table}

\begin{table}[htbp]
  \centering
  \small
  \begin{adjustbox}{max width=\textwidth}
  \gentable{tab_v2_e3a_cubebox_v1}
  \end{adjustbox}
  \caption{E3a on cube in box v1, as Table~\ref{tab:v2e3asort}.}\label{tab:v2e3acube}
\end{table}

\subsection{The stack on the bench (E3b)}\label{sec:v2e3b}

Figure~\ref{fig:v2stack} and Tables~\ref{tab:v2e3bsort} and~\ref{tab:v2e3bcube} compare the stack with \drtc{} alone
under the mix, and
Tables~\ref{tab:v2e3bablsort} and~\ref{tab:v2e3bablcube} remove one part at a time. \smolvla{} and \groot{} ran the
stack as chosen, with RTC; the five policies without RTC ran it with temporal ensembling as the seam method.

\paragraph{The stack against \drtc{} alone.} The stack lowered the jerk for every policy on both tasks, to between
0.05 and 0.31 of \drtc{} alone (3 to 20 times smoother), and the largest single step by a factor of 2 to 6. It cut the
stalls of the slow policies: on sorting Diffusion Policy from 50 to 27\%, MolmoAct~2 from 57 to 32\%, \groot{} from 65
to 41\% and GR00T~N1.7 from 68 to 49\%, with the same pattern on cube in box. Recovery after the faults was as fast or
faster for five of the seven policies. Three costs remain. The arm deviated from the policy's plan by 0.7 to
2.3\degree, more than any single method, because the filter, ensembling and the ramp each add their share. \smolvla's
stalls doubled (15 to 33\% on sorting, 14 to 31\% on cube in box). And \pizerofive{} recovered more slowly on sorting
(1.0 to 2.2~s), with more stalls there (40 to 46\%) but fewer on cube in box (40 to 33\%).

\paragraph{The parts.} Adaptive playback is what removes the stalls: without it the stall fraction rose by 11 to 19
percentage points (median per group) in 12 of the 14 policy-task pairs. The Butterworth filter is what removes most of
the jerk: without it the jerk was 1.7 to 2.4 times higher, in 14 of 14 pairs. The resume ramp mattered little in the
stack (median change of the jerk $-4$ to $+22\%$), because adaptive playback already removes most of the restarts it
smooths. The seam method behaved differently by kind. Removing RTC \emph{lowered} the stalls on both tasks, of \smolvla{} by 16 to
20 points and of \groot{} by 3, at no cost in jerk (1.01 to 1.03 times); \smolvla{} without RTC stalled no more than with
\drtc{} alone, and on cube in box \groot{} recovered from the latency spike in 0.7~s without RTC against 2.3~s with it.
Removing ensembling changed little on sorting and raised the jerk by 1.28 times on cube in box.

\paragraph{What this means for the stack.} By its pre-registered rule the stack includes RTC, but E3b shows that RTC is
the one part that does not earn its place on this setup: it makes each request slower, and with a round trip
dominated by the network the extra delay costs more stalls than the smoother seams save. The stack without RTC
(adaptive playback, temporal ensembling, the Butterworth filter and the resume ramp) stalled less than the chosen stack on all four RTC
policy-task pairs (\smolvla{} 16 to 20 points, \groot{} 3), at the same jerk. The rule is kept as it was fixed, and this result is reported as its correction: the robot experiment E4 uses the
stack without RTC for all four of its policies (Section~\ref{sec:v2robot}).

Without faults the stack still stalled \groot{} (20--23\%) and GR00T~N1.7 (36\%): the chunk of these policies is barely
longer than their round trip even on a good network (Section~\ref{sec:shortchunk}). The other policies stalled for
less than 13\% of a fault-free run, \smolvla{} excepted (22\%, with RTC).

\begin{figure}[htbp]
  \centering
  \genfigure{fig_v2_stack}
  \caption{E3b: each policy's median stall fraction under the fault mix with \drtc{} alone (grey) and with the stack
           (orange), and how many times lower the stack's jerk was (right).}
  \label{fig:v2stack}
\end{figure}

\begin{table}[htbp]
  \centering
  \small
  \begin{adjustbox}{max width=\textwidth}
  \gentable{tab_v2_e3b_sort_one}
  \end{adjustbox}
  \caption{E3b on sorting: the stack against \drtc{} alone under the fault mix, and the stack without faults (medians
           over 8 runs; jerk: the stack's median over \drtc's; recovery: median over the mix's three faults). Seam: the
           stack's seam method, RTC where the policy supports it and temporal ensembling otherwise.}\label{tab:v2e3bsort}
\end{table}

\begin{table}[htbp]
  \centering
  \small
  \begin{adjustbox}{max width=\textwidth}
  \gentable{tab_v2_e3b_cubebox_v1}
  \end{adjustbox}
  \caption{E3b on cube in box v1, as Table~\ref{tab:v2e3bsort}.}\label{tab:v2e3bcube}
\end{table}

\begin{table}[htbp]
  \centering
  \small
  \gentable{tab_v2_e3b_abl_sort_one}
  \caption{E3b ablations on sorting: the stack without one part, against the full stack under the mix (number of
           policies in which the jerk or the stalls were higher without the part; median ratio of the jerk, median
           change of the stall fraction and of the recovery time).}\label{tab:v2e3bablsort}
\end{table}

\begin{table}[htbp]
  \centering
  \small
  \gentable{tab_v2_e3b_abl_cubebox_v1}
  \caption{E3b ablations on cube in box v1, as Table~\ref{tab:v2e3bablsort}.}\label{tab:v2e3bablcube}
\end{table}

